\documentclass[UTF8]{ctexart}
\usepackage{xeCJK}
\usepackage{geometry}
\geometry{left=2cm,right=2cm,top=2.5cm,bottom=2.5cm}
\setlength{\headheight}{12.7pt}

\usepackage{titlesec}
\titleformat{\section}
  {\normalfont\large\bfseries}
  {\thesection}
  {1em}
  {}
\titlespacing*{\section}
  {0pt}
  {3.5ex plus 1ex minus .2ex}
  {2.3ex plus .2ex}

\usepackage{amsmath}
\usepackage{amssymb}
\usepackage{amsthm}
\usepackage{enumitem}
\setlist[enumerate]{itemsep=0pt,topsep=0pt,parsep=0pt,leftmargin=0pt,itemindent=2.5em}
\setlist[itemize]{itemsep=0pt,topsep=0pt,parsep=0pt,leftmargin=0pt,itemindent=2.5em}
\usepackage{fancyhdr}
\pagestyle{fancy}
\fancyhf{}
\usepackage{graphicx}
\usepackage{hyperref}
\usepackage{indentfirst}
\usepackage{lmodern}


\begin{document}
\setlength{\abovedisplayskip}{1pt}
\setlength{\belowdisplayskip}{1pt}

\section{陪集与Lagrange定理}

\hypertarget{sec:定义1.1}{}
\noindent\textbf{定义1.1 (陪集)}

给定群$G$与\href{02 子群#nameddest=sec:定义1.1}{子群} $H$,
对任意的$g\in G$, 分别称
\[
    gH=\{gh\mid h\in H\},\quad Hg=\{hg\mid h\in H\}\\[1pt]
\]
为$H$的\textbf{左陪集}和\textbf{右陪集}. 显然陪集与原子群等势.

\vspace{10pt}

\hypertarget{sec:定义1.2}{}
\noindent\textbf{定义1.2}

给定群$G$与子群$H$, 定义商集
\[
    \,\!^{\displaystyle G}\!/_{\displaystyle H}
    \overset{\mathrm{def}}{=}\{gH\mid g\in G\},
\]
它恰好是$G$的一个划分.

\noindent{\kaishu 验证.}

\begin{enumerate}
    \item 显然任意的$gH$非空,
    \item 如果$g_1H\neq g_2H$, 再假设存在$g\in g_1H\cap g_2H$, 那么
    存在$h_1,h_2\in H$使得
    \[
        g=g_1h_1=g_2h_2\quad\implies\quad g_1H=g_2h_2h_1^{-1}H=g_2H,
    \]
    矛盾,
    \item 对任意的$g\in G$, 有$g\in gH$. \hfill $\qedsymbol$
\end{enumerate}

\vspace{15pt}

\hypertarget{sec:定理1.1}{}
\noindent\textbf{定理1.1 (Lagrange定理)}

任给群$G$与子群$H$, 有
$|G|=\left|\,\!^{\displaystyle G}\!/_{\displaystyle H}\right|\cdot|H|$.

\vspace{3pt}
\noindent{\kaishu 证明.}

\vspace{3pt}
取$\,\!^{\displaystyle G}\!/_{\displaystyle H}$上的选择函数$f$,
然后定义映射
\vspace{3pt}
\[
    \,\!^{\displaystyle G}\!/_{\displaystyle H}\times H\to G,\quad
    (P,h)\mapsto f(P)\cdot h,
\]
下证它是双射.

\begin{enumerate}
    \item 如果$f(P_1)\cdot h_1=f(P_2)\cdot h_2$,
    那么它同时属于$P_1$和$P_2$, 所以$P_1=P_2$,
    \vspace{3pt}
    \item 对任意的$g\in G$, 由$f(gH)=gH$, 存在$h\in H$使得
    $f(gH)=gh$, 从而$g=f(gH)\cdot h^{-1}$.
\end{enumerate}

\vspace{6pt}
所以$|G|=\left|\,\!^{\displaystyle G}\!/_{\displaystyle H}\times H\right|
=\left|\,\!^{\displaystyle G}\!/_{\displaystyle H}\right|\cdot|H|$.
\hfill $\qedsymbol$

\vspace{15pt}

\hypertarget{sec:定义1.3}{}
\noindent\textbf{定义1.3 (指标)}

给定群$G$与子群$H$, 定义
\[
    [G:H]\,\overset{\mathrm{def}}{=}\,
    \left|\,\!^{\displaystyle G}\!/_{\displaystyle H}\right|,
\]
称之为$H$在$G$中的\textbf{指标}.

显然当$G$有限时, $[G:H]=|G|\,/\,|H|$.

\end{document}